\subsection{Conical symplectic singularities}

The main class of examples of algebras $A$ we will be interested in is $A={\mathcal O}(X)$, where~$X$ is a (normal) conical symplectic singularity in the sense of Beauville (see~\cite{Ka, L2} and references therein for their definition and basic properties).

We will need the following lemma.

\begin{Lemma}[{\cite[Proposition~2.5]{L2}}]\label{inn} If $X$ is a conical symplectic singularity then every Poisson derivation of $A:=\mathcal{O}(X)$ is inner.
\end{Lemma}

\begin{proof} Here is another proof. Any such homogeneous derivation is defined by an (a priori multivalued) holomorphic Hamiltonian $H$ on the big symplectic leaf $X^\circ$ of $X$. Moreover, $H$ is, in fact, single-valued, since by~\cite{Na1} the algebraic fundamental group of $X^\circ$ is finite (indeed, we have a homomorphism $\phi\colon \pi_1(X^\circ,x_0)\to {\mathbb C}$ given by $\phi(\gamma)=\gamma(H)-H$, where $\gamma(H)$ is the analytic continuation of~$H$ along~$\gamma$, and $\phi=0$ since $\pi_1(X^\circ,x_0)_{\rm alg}$ is finite). Then by normality of~$X$ the function~$H$ extends to the whole~$X$, to a homogeneous holomorphic function on~$X$. This function is a holomorphic section of a line bundle on ${\mathbb P}X$, which is algebraic by the GAGA theorem, so $H\in \mathcal{O}(X)$, as claimed.
\end{proof}

It is shown in \cite{L2} that for conical symplectic singularities there is a very nice parametrization of filtered quantizations of $A$, which is the same as Namikawa's parametrization of filtered Poisson deformations~$X$~\cite{Na3}. Namely, they are parametrized by the finite dimensional space $\mathfrak{P}=H^2\big(\widetilde X^{\rm reg},{\mathbb C}\big)$, where $\widetilde X^{\rm reg}$ is the smooth part of the ${\mathbb Q}$-terminalization $\widetilde X$ of~$X$, modulo the action of a finite real reflection group~$W$ called the Namikawa Weyl group. Denote the quantization corresponding to $\lambda\in \mathfrak{P}$ by~${\mathbf A}_\lambda$. It is easy to show that ${\mathbf A}_\lambda$ is ${\mathbb Z}/2$-invariant.

Since $W$ is a reflection group, the space $\mathfrak{P}/W$ bijectively parametrizing quantizations is an affine space, whose coordinates are independent homogeneous generating invariants $p_1,\dots,p_r\in {\mathbb C}[\mathfrak{P}]^W$ of degrees $d_1,\dots,d_r$.

It is not hard to see that ${\mathbf A}_\lambda$ is even if and only if $\lambda$ and $-\lambda$ are equivalent with respect to the Namikawa Weyl group, i.e., if and only if there is an element $w\in W$ such that $w\lambda=-\lambda$ (see \cite{BPR}). This is equivalent to saying that $p_i=0$ whenever $d_i$ is odd, which defines a subspace in
$\mathfrak{P}/W$.

\begin{Remark}Lemma \ref{inn} implies that in the case of conical symplectic singularities the antiautomorphism $\sigma\colon {\mathbf A}\to {\mathbf A}$ such that $\sigma^2=s$ and $\operatorname{gr}\sigma={\rm i}^d$ is unique if exists. Indeed, if~$\sigma_1$,~$\sigma_2$ are two such antiautomorphisms then $g:=\sigma_1\circ \sigma_2^{-1}$ is an automorphism commuting with~$s$ such that $\operatorname{gr} g=1$. Hence $\log (g)$ is a linear combination of derivations of~${\mathbf A}$ of degree $\le -2$. If $g\ne 1$, then in the leading order $\log(g)$ yields a nonzero homogeneous derivation $D\colon A\to A$ of degree $\le -2$. By Lemma~\ref{inn}, it is inner, so given by a Hamiltonian of degree $\le 0$. But such a~Hamiltonian is necessarily constant, so $D=0$, a contradiction.

The same argument proves that if $g\colon {\mathbf A}\to {\mathbf A}$ is a filtration-preserving
automorphism commuting with $s$ then $g$ is completely determined by $\operatorname{gr}g$.
\end{Remark}
